\chapter{Kết luận}

Quá trình thực tập tại phòng thí nghiệm kiểm định chất lượng sản phẩm của Công ty TNHH Quốc Tế Khánh Sinh vận dụng trực tiếp các kiến thức Hóa phân tích vào môi trường sản xuất thực tế. Việc tham gia lấy mẫu, chuẩn bị mẫu, pha hóa chất, thực hiện các phép phân tích và xử lý số liệu là cơ hội tiếp cận các chỉ tiêu kiểm nghiệm quan trọng của phân bón như nitơ tổng số, phospho hữu hiệu, kali hữu hiệu, carbon hữu cơ, chất hữu cơ và độ ẩm. Đồng thời, quá trình thực tập còn giúp làm quen với việc sử dụng các thiết bị như cân phân tích, hệ thống chưng cất, máy quang kế ngọn lửa, tủ sấy và các dụng cụ phân tích thể tích.

Bên cạnh kiến thức chuyên môn, nhiều kỹ năng nghề nghiệp cũng được hình thành và củng cố, bao gồm thao tác phòng thí nghiệm chính xác, quản lý mẫu, ghi nhận số liệu, tính toán kết quả, đánh giá độ lặp và đối chiếu kết quả thành phẩm với chỉ tiêu đăng ký. Qua đó, mối liên hệ giữa lý thuyết Hóa phân tích với thực tế kiểm soát chất lượng, đặc biệt là vai trò của việc chuẩn hóa thao tác, kiểm soát sai số và đảm bảo khả năng truy xuất của kết quả được thể hiện rõ hơn.

Trong quá trình thực tập vẫn ghi nhận một số hạn chế như một số phép thử còn phụ thuộc nhiều vào thao tác thủ công, một số mẫu có độ lặp chưa đạt yêu cầu hoặc kết quả thành phẩm khác với mức công bố. Do thời gian thực tập còn hạn chế, các trường hợp này chưa được phân tích lại đầy đủ để xác định nguyên nhân chính xác. Vì vậy, một số giải pháp được đề xuất gồm chuẩn hóa thao tác, tăng cường kiểm tra thiết bị và đường chuẩn, sử dụng mẫu QC định kỳ, thực hiện lại các mẫu không đạt và từng bước số hóa việc ghi nhận, tính toán và lưu trữ dữ liệu. Nhìn chung, kỳ thực tập đã giúp củng cố nền tảng chuyên môn, nâng cao kỹ năng thực hành và hình thành tư duy kiểm soát chất lượng có hệ thống trong lĩnh vực phân tích phân bón.
